\section{Introduction}

Ce rapport dresse un bilan de l'approche Python historique presente dans \path{full_soft/}. Cette approche est analysee comme une logique de type Voiture Jaune, c'est-a-dire une architecture directe, non ROS, concue autour de scripts Python, d'interfaces materielles et d'une boucle de controle embarquee sur Raspberry Pi.

Le rapport s'interesse a son usage possible ou historique avec la Voiture Blue. Il ne remplace pas la documentation APEX actuelle: il vise a comprendre ce que contient \path{full_soft/}, comment cette pile fonctionne, ce qu'elle apporte, et quelles limites elle presente par rapport a la pile ROS 2 moderne.

Les faits decrits ici proviennent principalement des fichiers \path{full_soft/code/main.py}, \path{full_soft/code/main_refactored.py}, \path{full_soft/code/VehicleAlgorithm_v1.py}, du dossier \path{full_soft/code/algorithm}, du dossier \path{full_soft/code/Interfaces}, des fichiers de configuration et de la documentation materielle presente sous \path{full_soft/docs} et \path{full_soft/documentation}.

